\begin{figure}[H]
\centering
\begin{tikzpicture}[
  >=Latex,
  every node/.style={font=\small},
  box/.style={
    draw,
    rounded corners=2pt,
    align=center,
    inner sep=6pt,
    text width=5.45cm,
    minimum height=1.38cm
  },
  narrow/.style={
    draw,
    rounded corners=2pt,
    align=center,
    inner sep=6pt,
    text width=4.15cm,
    minimum height=1.15cm
  },
  flow/.style={->,thick}
]

\node[box] (true) at (2.9,1.45) {
\textbf{精确反向后验均值}\\[3pt]
$\displaystyle
\tilde{\boldsymbol{\mu}}_t
=
\frac{1}{\sqrt{\alpha_t}}
\left(
\mathbf{x}_t-
\frac{\beta_t}{\sqrt{1-\bar{\alpha}_t}}
\,\boldsymbol{\epsilon}
\right)$
};

\node[box] (learned) at (9.1,1.45) {
\textbf{学习到的反向均值}\\[3pt]
$\displaystyle
\boldsymbol{\mu}_{\theta}
=
\frac{1}{\sqrt{\alpha_t}}
\left(
\mathbf{x}_t-
\frac{\beta_t}{\sqrt{1-\bar{\alpha}_t}}
\,\boldsymbol{\epsilon}_{\theta}(\mathbf{x}_t,t)
\right)$
};

\draw[flow] (true) -- (learned);

\node[narrow] (samevar) at (2.9,-0.95) {
\textbf{共享的固定方差}\\[2pt]
$q:\ \tilde{\beta}_t\mathbf I
\qquad
p_\theta:\ \tilde{\beta}_t\mathbf I$
};

\node[narrow] (meanmse) at (9.1,-0.95) {
\textbf{高斯 KL 化为\\均值匹配}\\[2pt]
$\displaystyle
\mathrm{KL}(q\,\|\,p_\theta)
\propto
\|\tilde{\boldsymbol{\mu}}_t-\boldsymbol{\mu}_\theta\|^2$
};

\draw[flow] (true.south) -- (samevar.north);
\draw[flow] (learned.south) -- (meanmse.north);
\draw[flow] (samevar) -- (meanmse);

\node[box] (noise) at (6.0,-3.45) {
\textbf{共享系数只留下时间依赖的缩放}\\[3pt]
$\displaystyle
\|\tilde{\boldsymbol{\mu}}_t-\boldsymbol{\mu}_\theta\|^2
\propto
\|\boldsymbol{\epsilon}
-\boldsymbol{\epsilon}_{\theta}(\mathbf{x}_t,t)\|^2$
};

\draw[flow] (meanmse) -- (noise);

\node[narrow] (simple) at (6.0,-5.95) {
\textbf{实用 DDPM 目标}\\[3pt]
$\displaystyle
\mathcal L_{\mathrm{simple}}
=
\mathbb E
\left[
\|\boldsymbol{\epsilon}
-\boldsymbol{\epsilon}_{\theta}(\mathbf{x}_t,t)\|^2
\right]$
};

\draw[flow] (noise) -- (simple);

\end{tikzpicture}

\caption{\textbf{从精确后验匹配到噪声预测。}
把精确反向均值与学习到的反向均值写成同一形式，并使用相同的固定方差，局部高斯 KL 就转化为均值平方误差匹配。由于两者共享系数，该目标进一步化为真实噪声 $\boldsymbol{\epsilon}$ 与其预测 $\boldsymbol{\epsilon}_\theta(\mathbf{x}_t,t)$ 之间带时间权重的平方误差。常用的 $\mathcal{L}_{\mathrm{simple}}$ 目标会去掉这一时间依赖权重。}
\label{fig:ddpm-noise-prediction-bridge}
\end{figure}